% Generated draft: requires mathematical and source review.
\documentclass[11pt]{article}
\usepackage[margin=25mm]{geometry}
\usepackage{fontspec}
\setmainfont{Latin Modern Roman}
\usepackage{amsmath,amssymb,mathtools,unicode-math}
\setmathfont{Latin Modern Math}
\usepackage{enumitem,needspace,xurl,hyperref}
\setlength{\parindent}{0pt}
\setlength{\parskip}{5pt}
\setlength{\emergencystretch}{3em}
\newcommand{\N}{\mathbb{N}}
\newcommand{\Z}{\mathbb{Z}}
\newcommand{\Q}{\mathbb{Q}}
\newcommand{\R}{\mathbb{R}}
\newcommand{\C}{\mathbb{C}}
\newcommand{\F}{\mathbb{F}}
\newcommand{\A}{\mathbb{A}}
\newcommand{\PP}{\mathbb{P}}
\newcommand{\E}{\mathbb{E}}
\newcommand{\eps}{\varepsilon}
\newcommand{\dd}{\,\mathrm{d}}
\begin{document}
\section*{UnsolvedMath ID: 30004215}
\textbf{Source problem code:} catalogue entry
\subsection{Definitions and mathematical statement}
For integers \(1\leq k\leq n\), let \(\operatorname{Gr}(k,n)(\R)\) denote the Grassmannian of \(k\)-dimensional linear subspaces of \(\R^n\). Its totally nonnegative part is
\[
\operatorname{Gr}_{\geq 0}(k,n)=\{V\in\operatorname{Gr}(k,n)(\R):\text{the Pl\"ucker coordinates of }V\text{ admit a common representative with all coordinates nonnegative}\}.
\]
Strata of this set can be specified by requiring designated Pl\"ucker coordinates to vanish and the remaining coordinates to be positive. For a linear map \(Z:\R^n\to\R^r\), the induced map
\[
V\longmapsto Z(V)\in\operatorname{Gr}(k,r)(\R)
\]
is defined only where \(\dim Z(V)=k\). Images of totally nonnegative Grassmannians or their strata under such maps provide the relevant Grassmann-polytope constructions. The supplied record does not specify the permitted parameters, strata, hypotheses on \(Z\), or image and closure conventions; these must be recovered from the cited sources before fixing the intended class.

A positive geometry consists of a complex projective algebraic variety \(X\), defined over \(\R\), together with an appropriate oriented real semialgebraic region \(P\subseteq X(\R)\). Its defining canonical-form requirement is the existence of a unique nonzero rational differential form \(\Omega(X,P)\) of degree \(\dim X\), with only logarithmic poles along the algebraic boundary and no poles elsewhere, whose residues on boundary components are the canonical forms of the induced lower-dimensional positive geometries. The recursion ends at oriented points, with canonical form \(+1\) or \(-1\). Applying this definition requires specification of the ambient variety, orientation, and boundary geometries, not merely the underlying point set.

\textbf{Question as supplied.} Is every Grassmann polytope a positive geometry, including the amplituhedra singled out in the supplied background? Determine the topology and combinatorics of Grassmann polytopes in a form analogous to the theory of ordinary convex polytopes.

More explicitly, if \(\mathcal G\) denotes the intended class of Grassmann polytopes with their intended ambient varieties \(X_P\), the first question asks whether
\[
\forall P\in\mathcal G,\qquad (X_P,P)\text{ is a positive geometry}.
\]
This is a formulation template: the source record does not determine \(\mathcal G\), \(X_P\), or all required boundary conventions precisely enough to turn it into a fully specified universal assertion. The topology and combinatorics request is likewise a research objective; the record does not specify particular invariants, a face relation, or a proposed classification.
\subsection{Short English statement}
The supplied problem asks whether every Grassmann polytope has the recursive canonical-form structure of a positive geometry, and seeks a theory of its topology and combinatorics analogous to convex-polytope theory. The intended construction and boundary conventions need clarification. The supplied curation labels the general assertion open; this draft does not independently verify that status.
\subsection{Sources}
\begin{itemize}
\item Supplied local UnsolvedMath registry record, permanent ID \(30004215\), title \emph{Grassmann Polytopes as Positive Geometries}, source problem code \texttt{catalogue entry}, in the collection \emph{Oberwolfach Reports Open Problems}. The record gives proposed year 2019 and does not name a proposer.
\item Supplied report citation: \emph{Geometric, Algebraic, and Topological Combinatorics} (2020), \emph{Oberwolfach Reports}, pp.\ 2395--2472, DOI \texttt{10.4171/owr/2019/39}. \url{https://doi.org/10.4171/owr/2019/39}. The literature-source list labels this reference ``Oberwolfach Report 39, Grassmann polytopes.''
\item Supplied literature reference: Arkani-Hamed, Bai, and Lam, positive geometries. \url{https://doi.org/10.1007/JHEP11%282017%29039}.
\item Literature assessment checked at \texttt{2026-08-22}: the supplied dated curation labels the general positive-geometry assertion open, particularly for amplituhedra, and states that special cases and canonical-form results do not establish the general claim. This is supplied dated curation, not an independently verified literature assessment; the record explicitly requests independent verification of time-sensitive status and cited claims.
\item The registry labels the statement \texttt{corrected\_verified}, explaining that mathematical content was preserved while notation, typography, and sentence structure were repaired. That supplied editorial label does not constitute mathematical verification of this draft.
\end{itemize}
\end{document}
